\documentclass[11pt]{article}
\usepackage{amsmath,amsfonts}
\usepackage{graphicx}
%\usepackage{xcolor}
%\definecolor{gray}{rgb}{.75,.75,.75}
%\usepackage[landscape]{geometry}
%\usepackage{hyperref}
%\usepackage[bookmarks=true, pdfborder={0 0 .5}, urlbordercolor=gray]{hyperref}

\renewcommand{\thefootnote}{\fnsymbol{footnote}}


\addtolength{\topmargin}{-.6in}
\addtolength{\textheight}{1.0in}
\addtolength{\oddsidemargin}{-.35in}
\addtolength{\textwidth}{.7in}
\pagestyle{empty}

\newcommand{\ds}{\displaystyle}
\newcommand{\ts}{\textstyle}

\newcommand{\Z}{\mathbb{Z}}
\newcommand{\R}{\mathbb{R}}
\newcommand{\C}{\mathbb{C}}
\newcommand{\EE}{\mathcal{E}}
\newcommand{\tZ}{\tilde{\Z}}

\newcommand{\Sym}{\mathrm{Sym}}

\renewcommand{\circle}[1]{{\bigcirc\hspace{-.75em}#1}}

\begin{document}

\footnote[0]{Notes by Peter Magyar \ {\tt magyar@math.msu.edu}}
\vspace{-1em}

\centerline{\bf Math 133\hspace{1.5em} \hfill{\Large\bf
Improper Integrals}\hfill Stewart \S7.8}
\vspace{1em}

\noindent 
{\bf Integrals near a vertical asymptote.} What happens if we take the integral of a function over an interval containing a vertical asymptote, such as:
$$
I \ = \ \int_0^2 \frac{1}{x}\,dx \ =\ ??
$$
Algebraically, we would get $I=\ln|2| -\ln|0|$, but $\ln(0)$ is undefined.
Numerically, the Riemann sum for $I$ does not converge, because of the very large values of $f(x)$ near $x=0$. Geometrically, $I$ measures a
region (the positive area in the graph on the next page) which stretches infinitely 
along the asymptote $x=0$, and the meaning of such an
infinitely extended area is not clear.
 
Our previous definitions fail to give meaning to this integral, so we give a new definition:
$$
\int_0^2 \frac{1}{x}\,dx \ =\  \lim_{r\to 0^+} \int_r^2 \frac{1}{x}\,dx.
$$
That is, we take the integral over the interval $x\in [r,2]$ where the function is continuous, then take the limit as $r$ squeezes up against the asymptote $x=0$ from the right.
Now, $\int_r^2\frac1x\,dx= \ln|2|-\ln|r|$, 
and $\lim_{r\to 0^+} \ln(r)=-\infty$, meaning $\ln(r)$ becomes a larger and larger
negative number, so the improper integral is:\footnote{
We take $\ln(2)$ minus a larger and larger negative number; 
and this equals a larger and larger positive number, denoted by $\infty$.
}
$$
\int_0^2 \frac1x\,dx\ =\
\lim_{r\to 0^+} \ln(2)-\ln(r) 
\ =\ \infty.
$$

This says that the total area under the graph $y=\frac1x$ and above $[0,2]$ is infinite: no matter how many square units of paint are put on this region, there will still be unpainted area high up next to the asymptote.
\\[1em]
{\bf General definition:}
If the function $f(x)$ has a vertical asymptote near $x=q$,
we define the {\it improper integral of vertical type}:
\begin{itemize}
\item on an interval $[a,q]$ as \ \
$\ds
\int_a^q f(x)\,dx \ =\ \lim_{r\to q^-}\int_a^r f(x)\,dx;
$
\item on an interval $[q,b]$ as \ \
$\ds
\int_q^b f(x)\,dx \ =\ \lim_{r\to q^+}\int_r^b f(x)\,dx.
$
\item on an interval with $q\in(a,b)$ as \ \
$\ds
\int_a^b f(x)\,dx \ =\ \int_a^q f(x)\,dx + \int_q^b f(x)\,dx.
$
\end{itemize}
If such an integral has a finite value, we say it {\it converges};
if it is infinite or undefined, we say it {\it diverges}.
%
\newpage

\noindent
{\sc example:} Evaluate $\int_{-1}^2\frac 1x\,dx$.
This attempts to measure two infinite regions: one above $[0,2]$ along the positive $y$-axis, and another below $[-1,0]$ along the negative $y$-axis. 
\\[1em]
\centerline{
\includegraphics[width=2in]
{7-8-Inf-Area.png}
}
\vspace{0em}

\noindent
The improper integral avoids the asymptote from both sides:
$$
\int_{-1}^2 \frac1x\,dx 
\ =\ 
\int_{-1}^0 \frac1x\,dx + 
\int_{0}^2 \frac1x\,dx 
\ =\ 
\lim_{r\to 0^-}\int_{-1}^r \frac1x\,dx 
\,+\, \lim_{r\to 0^+}\int_r^2\frac1x\,dx.
$$
But when we try to calculate this, we get:
$$
\int_{-1}^2 \frac1x\,dx \ =\ 
\left(\lim_{r\to0^-}  \ln|r|-\ln|\!-1| \right)
\,+\,
\left(\lim_{r\to 0^+} \ln|2|-\ln|r|\right)
\ =\ -\infty +\infty\,,
$$
which is an indeterminate form: the integral is truly undefined.
We have no clear meaning for an infinite positive area canceled by an infinite negative area.
In particular, the naive answer is wrong:
$$
\int_{-1}^2 \frac1x\,dx \ =\ \text{undefined} \ \neq\ \ln|2|-\ln|\!-1|.
$$
\\[.5em]
\noindent
{\sc example:} Evaluate $\int_{1}^2 \frac{1}{\sqrt{x-1}}\,dx$.
Since the vertical asymptote is $x=1$, we have:
$$
\int_{1}^2 \frac{1}{\sqrt{x-1}}\,dx \ =\ 
\lim_{r\to 1^+}\int_{r}^2 \frac{1}{\sqrt{x-1}}\,dx 
\ =\ 
\lim_{r\to 1^+} \left.2\sqrt{x{-}1}\,\right|_{x=r}^{x=2}
$$
$$
\ =\ 
\lim_{r\to 1^+} 2\sqrt{2{-}1}-2\sqrt{r{-}1}
\ =\ 2 - 0 \ =\ 2.
$$
In this case, the region has a finite area of 2, even though it stretches infinitely high along the vertical asymptote. Thus, if we have enough paint for 2 square units, and paint higher and higher parts of the region using less and less paint, we never run out.

\newpage 

\noindent
{\bf Integrals near a horizontal asymptote.}
If $y=f(x)$ has $y=0$ as a horizontal asymptote, we can define  
{\it improper integrals of horizontal type}.
\begin{itemize}
\item If $\lim_{x\to\infty} f(x) = 0$, we define the integral on an interval $[a,\infty)$ as:
$$
\int_a^\infty f(x)\,dx
\ =\
\lim_{r\to\infty} \int_a^r f(x)\,dx.
$$
\item If $\lim_{x\to\infty} f(x)=0$, we define the integral over an interval
$(-\infty,a]$ as:
$$
\int_{-\infty}^a f(x)\,dx
\ =\
\lim_{r\to-\infty} \int_r^a f(x)\,dx.
$$
\item If $\lim_{x\to\pm\infty} f(x)=0$, we define its integral
over the whole real line $(-\infty,\infty)$
by splitting at any finite value $x=a$:
$$
\int_{-\infty}^{\infty} f(x)\,dx
\ =\
\int_{-\infty}^a f(x)\,dx + \int_a^{\infty} f(x)\,dx \qquad \text{for any $a$.}
$$
\end{itemize}

\noindent
{\sc example:}
$
\int_1^\infty \!\frac1{x^2} \,dx
\ =\
\lim\limits_{r\to\infty} \int_1^r \frac{1}{x^2}\,dx
\ =\
\lim\limits_{r\to\infty} \left(-\left.\frac1{x}\right)\right|_{x=1}^{x=r}
\ =\ \lim_{r\to\infty} -\frac1r + \frac11 \ =\ 1\,.
$
\\[.2em]
This integral measures a region which stretches infinitely along the $x$-axis above $[1,\infty)$, but which has a finite total area of 1.

On the other hand $\int_1^\infty \frac{1}{\sqrt x}\,dx
\ =\ \lim_{r\to\infty} \left.2\sqrt x\ \right|_{x=1}^{x=r}\ =\infty$.
In fact, $\int_{1}^\infty \frac{1}{x^p}\,dx$ is finite if $p>1$, but is infinite if $p\leq 1$.
Informally, the faster $f(x)$ shrinks as $x\to\infty$, 
the easier it is for the integral to converge to a finite value. 
\\[1em]
{\sc example:}
$\int_0^\infty e^{-x}\,dx =\lim_{r\to\infty} \int_0^r e^{-x}\,dx
= \lim_{r\to\infty} \left.-e^{-x}\right|_{x=0}^{x=r} = -0-(-1)=1$.
\\[.2em]
It is not surprising that this converges, because $e^{-x}$ shrinks faster than $\frac1{x^p}$ for any $p$.
\\[1em]
{\sc example:}
$$
\int_{-\infty}^{\infty} \frac{1}{1+x^2}\,dx
\ =\
\lim_{r\to-\infty}\int_{r}^0 \frac{1}{1+x^2}\,dx
\ +\
\lim_{r\to\infty}\int_{0}^r \frac{1}{1+x^2}\,dx
$$
$$
\ =\
\lim_{r\to-\infty} \left.\tan^{-1}(x)\right|_{x=r}^{x=0}
\ +\ 
\lim_{r\to\infty} \left.\tan^{-1}(x)\right|_{x=0}^{x=r}
\  =\
\left(0-(-\tfrac{\pi}2)\right) + \left(\tfrac{\pi}2-0\right)
\ =\ \pi.
$$
Remarkably, the total area under $y=\frac 1{1+x^2}$ turns out to be $\pi$, same as a unit circle!
\\[1em]
\centerline{
\includegraphics[width=3in]
{7-8-Bell.png}
}

\newpage

\noindent
{\bf Comparison tests for convergence.}  
Sometimes an improper integral is too complicated to find an algebraic antiderivative, but we can still be sure it converges
because the infinite region measured fits inside a larger region of known finite area.

For example, the Gaussian bell-curve integral
$\int_1^\infty e^{-x^2}\,dx$ cannot be integrated by an antiderivative. However, for $x\geq 1$, we have $x^2\geq x$,
so $e^{-x^2}\leq e^{-x}$: that is, the curve $y=e^{-x^2}$ lies below $y=e^{-x}$:
\\[1em]
\centerline{
\includegraphics[width=2.5in]
{7-8-Gaussian.png}
}
\\[.5em]
We can easily evaluate the area below the upper curve,
which shows that the smaller area under the
lower curve is finite, i.e.~the improper integral converges:
$$
\int_1^\infty e^{-x^2}\,dx
\ <\
\int_1^\infty e^{-x}\,dx
\ =\ 0-(-e^{-1}) \ =\ \tfrac1e \ \approx\ 0.37\ .
$$
{\it Direct Comparison Test:} Consider an improper integral
$\int_a^b g(x)$, with $a$ or $b$ infinite.
\begin{itemize}
\item
If $|f(x)|\,{\leq}\, g(x)$ for $x\,{\in}\,[a,b]$, and $\int_a^b g(x)\,dx$ converges,
then $\int_a^b f(x)\,dx$ converges.
\item If $f(x)\,{\geq}\,g(x)\,{\geq}\,0$  for $x\,{\in}\,[a,b]$ and $\int_a^b g(x)\,dx$ diverges,
then $\int_a^b f(x)\,dx$ diverges.
\end{itemize}

\noindent
The proof uses the Domination Rule for ordinary integrals (\S4.2),
plus some complications with limits.
\\[.5em]
{\sc example:} Does $\int_0^\infty \frac{4\sin(x)+1}{e^{2x}+x^2}\,dx$ converge? 
This function shrinks rapidly, since the top does not grow,
and the bottom grows exponentially; thus we guess that
the integral converges.
To prove this using the first part of the Test, we should 
bound \smash{$f(x)=\frac{4\sin(x)+1}{e^{2x}+x^2}$} 
inside the graph of a fairly simple comparison function
$g(x)=\frac{g_1(x)}{g_2(x)}$ with $|f(x)|\leq g(x)$.
Now, increasing the numerator of $f(x)$ and 
decreasing its denominator gives a larger fraction, so we take:
$$
\left|\frac{4\sin(x)+1}{e^{2x}+x^2}\right|
\ \leq\ 
\frac{5}{e^{2x}} \ =\  5e^{-2x}\,.
$$
The comparison integral converges: $\int_0^\infty 5e^{-2x}\,dx = \lim_{r\to\infty}
\left.\left(-\frac52 e^{-2x}\right)\right|_{x=0}^{x=r} = \frac52$; hence 
the original integral also converges: 
$$
\left|\int_0^\infty \frac{4\sin(x)+1}{e^{2x}+x^2}\,dx\right|
\ \leq\ \tfrac52\,.
$$
\\[1em]
By contrast, to prove {\it divergence} of a fractional $f(x)$, 
we would bound $f(x)\geq g(x)$, above a floor function $g(x)$ with
smaller numerator and larger denominator, and $\int g(x)\,dx=\infty$.
\newpage

\noindent
{\it Limit Comparison Test} or {\it Ratio Comparison Test:}
For functions $f(x), g(x)$ with $\lim\limits_{x\to\infty} \frac{f(x)}{g(x)} = L\neq 0$,
the improper integral $\int_a^\infty f(x)\,dx$ converges if and only if 
$\int_a^\infty g(x)\,dx$ converges.
\\[1em]
In the case that $g(x)\geq 0$, this is simply because, given $\lim_{x\to\infty} \frac{f(x)}{g(x)} = L$, we can take $x$ large enough so that 
$\frac12 L g(x)\leq f(x)\leq \frac32L g(x)$, and we can apply the Direct Comparison Test.

To apply this Test to $\int_a^\infty f(x)\,dx$ for a fraction $f(x)=\frac{f_1(x)}{f_2(x)}$,
we generally choose the comparison function
$g(x)=\frac{g_1(x)}{g_2(x)}$ where $g_1(x)$ is the
 largest term in $f_1(x)$, and likewise with $g_2(x)$ and
$f_2(x)$.
For example, for:
$$
f(x)=\frac{x^2-e^{-x}+\sin(x)}
{\sqrt{x^5+7}}\qquad\text{take}\qquad
g(x)=\frac{x^2}{\sqrt{x^5}}=x^{-1/2}\,,
$$
and we easily see $\lim_{x\to\infty} \frac{f(x)}{g(x)} = 1$ (see \S6.8).
We previously showed $\int_a^\infty x^{-1/2}\,dx$ diverges, so the original integral
$\int_a^\infty f(x)\,dx$ also diverges.

\end{document}
%%%%%%%%%%%  Picture  %%%%%%%%%%
\\[0em]
\centerline{
%\includegraphics[width=2in]
{1-8-Discont-iv.png}
\qquad \qquad 
%\includegraphics[width=2in]
{1-8-Discont-v.png} 
}
\vspace{0em}

\noindent
%%%%%%%%%%%%%%%%%%%%%%%%%

%%%%%%%%%%%  Table  %%%%%%%%%%
In fact, we get the following derivatives:
$$\begin{array}{|c||c|c|c|c|c|c|}
\hline &&&&&& \\[-.9em] 
f(x) & \sin(x) & \cos(x) & \tan(x)  & \sec(x) & \csc(x) & \cot(x)
\\[.2em] \hline &&&&&& \\[-.9em]
f'(x) &\cos(x) & -\sin(x) & \sec^2(x) & \tan(x)\sec(x) & -\cot(x)\csc(x) & -\csc^2(x)
\\[.1em] \hline
\end{array}$$
\\[1em]
%%%%%%%%%%%%%%%%%%%%%%%%%%



















